% Capítulo I
\chapter{O que são algarismos?}
\label{oque_sao_algarismos}
% Conteúdo do primeiro capítulo
\lettrine{E}{ste capítulo nos dará as ferramentas iniciais} e o vocabulário para discutirmos melhor alguns conceitos matemáticos. A discussão mais aprofundada do que são números e conjuntos numéricos ficará outra apostila. Enquanto isso, vamos relembrar alguns conceitos importantes do sistema numérico que usamos no nosso dia a dia.

\section{Sistema de numeração decimal}

Esse conteúdo é importantíssimo da mesma maneira que é indispensável primeiro aprender o alfabeto, e só depois aprender a ler e a escrever. 

O sistema numérico que usamos no nosso cotidiano é o \textbf{decimal}, ou seja, todos os números que conhecemos podem ser formados utilizando somente \textbf{10 números}: {0, 1, 2, 3, 4, 5, 6, 7, 8, 9}, que são chamados de “algarismos arábicos”; e \textbf{vírgula} (como veremos no capítulo de decimais). Em outra apostila de matemática mais completa, poderemos fazer uma discussão sobre o motivo de se utilizar o sistema decimal, e quais outros sistemas numéricos existem/existiram ao longo da história, assim como a origem dos números.

Porém, no momento precisamos entender o princípio básico de formação de números maiores do que 9, que envolvem a combinação de um ou mais algarismos, e como arranjá-los. Por exemplo, o número 10 é formado pela junção dos algarismos “1” e “0”; em 156 usamos “1”, “5” e “6”, e assim em diante. 

Formalizando um pouco essa noção de “formação” e “arranjo” dos números, podemos falar do conceito de ordem dos números, sendo eles: \textbf{unidades} (U), \textbf{dezenas} (D) e \textbf{centenas} (C) - que serão muito importantes para os capítulos seguintes. Não estamos interessados em definições formais e rigorosas no momento, em outro volume veremos essas definições e explicações mais aprofundadas. De modo geral, temos a seguinte tabela:


\begin{table}[H]
\begin{tabular}{|ccc|ccc|ccc|}
\hline
\multicolumn{3}{|c|}{3ª classe}                               & \multicolumn{3}{c|}{2ª classe}                               & \multicolumn{3}{c|}{1ª classe}                               \\ \hline
\multicolumn{3}{|c|}{milhões}                                 & \multicolumn{3}{c|}{milhares}                                & \multicolumn{3}{c|}{-}                                       \\ \hline
\multicolumn{1}{|c|}{9ª O} & \multicolumn{1}{c|}{8ª O} & 7ª O & \multicolumn{1}{c|}{6ª O} & \multicolumn{1}{c|}{5ª O} & 4ª O & \multicolumn{1}{c|}{3ª O} & \multicolumn{1}{c|}{2ª O} & 1ª O \\ \hline
\multicolumn{1}{|c|}{C}    & \multicolumn{1}{c|}{D}    & U    & \multicolumn{1}{c|}{C}    & \multicolumn{1}{c|}{D}    & U    & \multicolumn{1}{c|}{C}    & \multicolumn{1}{c|}{D}    & U    \\ \hline
\end{tabular}

\caption{Classe e Ordem (O) dos números decimais.}
\label{ordemclasse}
\end{table}

A Tabela \ref{ordemclasse} pode parecer um pouco confusa e desnecessária de início, mas vamos tentar entender melhor o que esse sistema tem a oferecer. Primeiro, temos que definir o que significa cada ordem:

\begin{enumerate}

    \item \textbf{Unidades (U)} podem ser entendidas como o valor “único” de um número, ou seja, formado por apenas um algarismo arábico (0 a 9). Por exemplo: 1 maçã, 5 pedras, 9 multas de trânsito etc. São representadas na nossa tabela na posição de 1ª ordem (veremos exemplos mais à frente para esclarecer).
    \item \textbf{Dezenas (D)} são compostas de 10 unidades (devido ao nosso sistema decimal), ou seja, a cada 10 unidades temos 1 dezena. Por exemplo: 2 dezenas representam 20 unidades; 5 dezenas representam 50 unidades  etc. 
    
    Vamos pegar um exemplo mais cotidiano (ou nem tanto). Caso você vá até o mercadinho perto de sua casa e peça 1 dezena de ovos, o Seu Manuel vai entender que você deseja 10 ovos, faz sentido? Sabemos que é mais comum pedirmos “1 dúzia” de ovos (que representam 12 ovos) ou “meia dúzia” de alguma coisa (que representa 6) mas, aqui, somos pessoas estranhas que compram 10 ovos.
    
    \item \textbf{Centenas (C)}, seguindo a mesma lógica, representam 100 unidades (ou 10 dezenas). Por exemplo: 2 centenas representam 200 unidades (ou 20 dezenas); 5 centenas representam 500 unidades (ou 50 dezenas) etc. Um exemplo clássico é quando vamos comprar salgadinhos e pedimos “1 cento”, que representam 100 salgadinhos - nesse caso abreviamos “centena” para “cento”, mas é a mesma lógica.

\end{enumerate}

Agora vamos pegar exemplos mais interessantes, ``misturando'' unidades, centenas e dezenas. Vejamos:

\begin{itemize}
    \item 79 = 70 + 9 (7 dezenas e 9 unidades)
    \item 156 = 100 + 56 = 100 + 50 + 6 (1 centena, 5 dezenas e 6 unidades)
    \item 678 = 600 + 78 = 600 + 70 + 8 (6 centenas, 7 dezenas e 8 unidades)
\end{itemize}

Caso fossemos utilizar a nossa tabela, ficaria da seguinte maneira: \\

\textbf{Exemplo.:} 49 = 40 + 9 = 4 dezenas e 9 unidades.

\begin{table}[H]
\begin{tabular}{|ccc|ccc|ccc|}
\hline
\multicolumn{3}{|c|}{3ª classe} & \multicolumn{3}{c|}{2ª classe}                                                                       & \multicolumn{3}{c|}{1ª classe}                                                                        \\ \hline
\multicolumn{3}{|c|}{milhões}                                                                         & \multicolumn{3}{c|}{milhares}                                                                        & \multicolumn{3}{c|}{-}                                                                                \\ \hline
\multicolumn{1}{|c|}{9ª O}                     & \multicolumn{1}{c|}{8ª O}                     & 7ª O & \multicolumn{1}{c|}{6ª O}                     & \multicolumn{1}{c|}{5ª O}                     & 4ª O & \multicolumn{1}{c|}{3ª O}                     & \multicolumn{1}{c|}{2ª O}                      & 1ª O \\ \hline
\multicolumn{1}{|c|}{C}                        & \multicolumn{1}{c|}{D}                        & U    & \multicolumn{1}{c|}{C}                        & \multicolumn{1}{c|}{D}                        & U    & \multicolumn{1}{c|}{C}                        & \multicolumn{1}{c|}{D}                         & U    \\ \hline
\rowcolor[HTML]{FFFC9E} 
\multicolumn{1}{|c|}{\cellcolor[HTML]{FFFC9E}} & \multicolumn{1}{c|}{\cellcolor[HTML]{FFFC9E}} &      & \multicolumn{1}{c|}{\cellcolor[HTML]{FFFC9E}} & \multicolumn{1}{c|}{\cellcolor[HTML]{FFFC9E}} &      & \multicolumn{1}{c|}{\cellcolor[HTML]{FFFC9E}} & \multicolumn{1}{c|}{\cellcolor[HTML]{FFFC9E}\textbf{4}} & \textbf{9}    \\ \hline
\end{tabular}
\end{table}

\textbf{Exemplo.:} 253 = 200 + 50 + 3 = 2 centenas, 5 dezenas e 3 unidades.

\begin{table}[H]
\begin{tabular}{|ccc|ccc|ccc|}
\hline
\multicolumn{3}{|c|}{3ª classe} & \multicolumn{3}{c|}{2ª classe} & \multicolumn{3}{c|}{1ª classe} \\ \hline
\multicolumn{3}{|c|}{milhões} & \multicolumn{3}{c|}{milhares} & \multicolumn{3}{c|}{-} \\ \hline
\multicolumn{1}{|c|}{9ª O} & \multicolumn{1}{c|}{8ª O} & 7ª O & \multicolumn{1}{c|}{6ª O} & \multicolumn{1}{c|}{5ª O} & 4ª O & \multicolumn{1}{c|}{3ª O} & \multicolumn{1}{c|}{2ª O} & 1ª O \\ \hline
\multicolumn{1}{|c|}{C} & \multicolumn{1}{c|}{D} & U & \multicolumn{1}{c|}{C} & \multicolumn{1}{c|}{D} & U & \multicolumn{1}{c|}{C} & \multicolumn{1}{c|}{D} & U \\ \hline
\rowcolor[HTML]{FFFC9E} 
\multicolumn{1}{|c|}{\cellcolor[HTML]{FFFC9E}} & \multicolumn{1}{c|}{\cellcolor[HTML]{FFFC9E}} &  & \multicolumn{1}{c|}{\cellcolor[HTML]{FFFC9E}} & \multicolumn{1}{c|}{\cellcolor[HTML]{FFFC9E}} &  & \multicolumn{1}{c|}{\cellcolor[HTML]{FFFC9E}\textbf{2}} & \multicolumn{1}{c|}{\cellcolor[HTML]{FFFC9E}\textbf{5}} & \textbf{3} \\ \hline
\end{tabular}
\end{table}

Perceba que podemos ter exemplos onde uma ou mais ordens podem ter o número zero, como nos casos:

\begin{itemize}
    \item 50 (5 dezenas e 0 unidades)
    \item 120 (1 centena, 2 dezenas e 0 unidades)
    \item 107 (1 centena, 0 dezenas e 7 unidades)
\end{itemize}

Assim temos vários outros exemplos e infinitas possibilidades.

Então, você deve estar se perguntando: \textit{para que servem as 2ª e 3ª classes}? Antes de continuar a leitura, tente adivinhar por conta própria.

Depois de muito esforço mental e dedicação, certamente, vamos apresentar o que são as outras classes. Sabemos que existem números maiores que 999 e caso tivéssemos somente a 1ª classe só poderíamos representar até essa grandeza (9 centenas, 9 dezenas e 9 unidades). Então, precisamos de classes maiores, que são as seguintes:

\subsection{Classe dos Milhares}

Apesar de termos mudado de classe, perceberemos que as ordens continuam sendo as mesmas. Dessa forma temos: \textbf{unidade de milhar} (1.000), \textbf{dezena de milhar} (10.000) e \textbf{centena de milhar} (100.000). Vamos já apresentar alguns exemplos para deixar mais claro, vejamos:

\begin{itemize}
    \item 1.255 = 1.000 + 255 = 1000 + 200 + 50 + 5\\
    (1 unidade de milhar, 2 centenas, 5 dezenas e 5 unidades)
    \item 12.534 = 10.000 + 2534 = 10.000 + 2.000 + 500 + 30 + 4\\
    (1 dezena de milhar, 2 unidades de milhar, 5 centenas, 3 dezenas e 4 unidades)
    \item 456.789 = 400.000 + 50.000 + 6.000 + 700 + 80 + 9\\
    (4 centenas de milhar, 5 dezenas de milhar, 6 unidades de milhar, 7 centenas, 8 dezenas e 9 unidades)
\end{itemize}

Já podemos perceber que quanto maior o nosso número, mais trabalhoso será decompor ele em classes e ordens. Usualmente utilizar a tabela facilitará nosso serviço. Veja:\\

\textbf{Exemplo.:} 1.356 = 1.000 + 300 + 50 + 6

\begin{table}[H]
\begin{tabular}{|ccc|ccc|ccc|}
\hline
\multicolumn{3}{|c|}{3ª classe} & \multicolumn{3}{c|}{2ª classe} & \multicolumn{3}{c|}{1ª classe} \\ \hline
\multicolumn{3}{|c|}{milhões} & \multicolumn{3}{c|}{milhares} & \multicolumn{3}{c|}{-} \\ \hline
\multicolumn{1}{|c|}{9ª O} & \multicolumn{1}{c|}{8ª O} & 7ª O & \multicolumn{1}{c|}{6ª O} & \multicolumn{1}{c|}{5ª O} & 4ª O & \multicolumn{1}{c|}{3ª O} & \multicolumn{1}{c|}{2ª O} & 1ª O \\ \hline
\multicolumn{1}{|c|}{C} & \multicolumn{1}{c|}{D} & U & \multicolumn{1}{c|}{C} & \multicolumn{1}{c|}{D} & U & \multicolumn{1}{c|}{C} & \multicolumn{1}{c|}{D} & U \\ \hline
\rowcolor[HTML]{FFFC9E} 
\multicolumn{1}{|c|}{\cellcolor[HTML]{FFFC9E}} & \multicolumn{1}{c|}{\cellcolor[HTML]{FFFC9E}} &  & \multicolumn{1}{c|}{\cellcolor[HTML]{FFFC9E}} & \multicolumn{1}{c|}{\cellcolor[HTML]{FFFC9E}} & \textbf{1} & \multicolumn{1}{c|}{\cellcolor[HTML]{FFFC9E}\textbf{3}} & \multicolumn{1}{c|}{\cellcolor[HTML]{FFFC9E}\textbf{5}} & \textbf{6} \\ \hline
\end{tabular}
\end{table}

Da mesma forma, podemos ter as ordens vazias como visto anteriormente. Por exemplo:\\
	
\textbf{Exemplo.:} 100.562 = 100.000 + 500 + 60 + 2

\begin{table}[H]
\begin{tabular}{|ccc|ccc|ccc|}
\hline
\multicolumn{3}{|c|}{3ª classe} & \multicolumn{3}{c|}{2ª classe} & \multicolumn{3}{c|}{1ª classe} \\ \hline
\multicolumn{3}{|c|}{milhões} & \multicolumn{3}{c|}{milhares} & \multicolumn{3}{c|}{-} \\ \hline
\multicolumn{1}{|c|}{9ª O} & \multicolumn{1}{c|}{8ª O} & 7ª O & \multicolumn{1}{c|}{6ª O} & \multicolumn{1}{c|}{5ª O} & 4ª O & \multicolumn{1}{c|}{3ª O} & \multicolumn{1}{c|}{2ª O} & 1ª O \\ \hline
\multicolumn{1}{|c|}{C} & \multicolumn{1}{c|}{D} & U & \multicolumn{1}{c|}{C} & \multicolumn{1}{c|}{D} & U & \multicolumn{1}{c|}{C} & \multicolumn{1}{c|}{D} & U \\ \hline
\rowcolor[HTML]{FFFC9E} 
\multicolumn{1}{|c|}{\cellcolor[HTML]{FFFC9E}\textbf{}} & \multicolumn{1}{c|}{\cellcolor[HTML]{FFFC9E}\textbf{}} & \textbf{} & \multicolumn{1}{c|}{\cellcolor[HTML]{FFFC9E}\textbf{1}} & \multicolumn{1}{c|}{\cellcolor[HTML]{FFFC9E}\textbf{0}} & \textbf{0} & \multicolumn{1}{c|}{\cellcolor[HTML]{FFFC9E}\textbf{5}} & \multicolumn{1}{c|}{\cellcolor[HTML]{FFFC9E}\textbf{6}} & \textbf{2} \\ \hline
\end{tabular}
\end{table}

Até agora conseguimos representar números até 999.999, mas também sabemos que existem números muito maiores do que esse. A partir disso, avançamos para a próxima classe: milhões.

\subsection{Classe dos Milhões}

Na classe dos milhões teremos o mesmo princípio que o da classe dos milhares, sendo: \textbf{unidade de milhão} (1.000.000), \textbf{dezena de milhão} (10.000.000) e \textbf{centena de milhão} (100.000.000). Essa é a classe na qual todos nós gostaríamos de poder representar nossa conta bancária, mas, para isso, ainda temos muita matemática pela frente. Vejamos alguns exemplos:

\begin{itemize}
    \item 1.256.893 = 1.000.000 + 200.000 + 50.000 + 6.000 + 800 + 90 + 3 \\ 
    (1 unidade de milhão, 2 centenas de milhar, 5 dezenas de milhar, 6 unidades de milhar, 8 centenas, 9 dezenas e 3 unidades)
\end{itemize}

Perceba que utilizar a tabela facilitará muito nosso trabalho, então vamos utilizá-la a partir de agora:

\begin{itemize}
    \item 3.875.263 = 3.000.000 + 800.000 + 70.000 + 5.000 + 200 + 60 + 3

\begin{table}[H]
\begin{tabular}{|ccc|ccc|ccc|}
\hline
\multicolumn{3}{|c|}{3ª classe} & \multicolumn{3}{c|}{2ª classe} & \multicolumn{3}{c|}{1ª classe} \\ \hline
\multicolumn{3}{|c|}{milhões} & \multicolumn{3}{c|}{milhares} & \multicolumn{3}{c|}{-} \\ \hline
\multicolumn{1}{|c|}{9ª O} & \multicolumn{1}{c|}{8ª O} & 7ª O & \multicolumn{1}{c|}{6ª O} & \multicolumn{1}{c|}{5ª O} & 4ª O & \multicolumn{1}{c|}{3ª O} & \multicolumn{1}{c|}{2ª O} & 1ª O \\ \hline
\multicolumn{1}{|c|}{C} & \multicolumn{1}{c|}{D} & U & \multicolumn{1}{c|}{C} & \multicolumn{1}{c|}{D} & U & \multicolumn{1}{c|}{C} & \multicolumn{1}{c|}{D} & U \\ \hline
\rowcolor[HTML]{FFFC9E} 
\multicolumn{1}{|c|}{\cellcolor[HTML]{FFFC9E}\textbf{}} & \multicolumn{1}{c|}{\cellcolor[HTML]{FFFC9E}\textbf{}} & \textbf{3} & \multicolumn{1}{c|}{\cellcolor[HTML]{FFFC9E}\textbf{8}} & \multicolumn{1}{c|}{\cellcolor[HTML]{FFFC9E}\textbf{7}} & \textbf{5} & \multicolumn{1}{c|}{\cellcolor[HTML]{FFFC9E}\textbf{2}} & \multicolumn{1}{c|}{\cellcolor[HTML]{FFFC9E}\textbf{6}} & \textbf{3} \\ \hline
\end{tabular}
\end{table}

    \item 45.874.674 = 40.000.000 + 5.000.000 + 800.000 + 70.000 + 4.000 + 600 + 70 + 4

\begin{table}[H]
\begin{tabular}{|ccc|ccc|ccc|}
\hline
\multicolumn{3}{|c|}{3ª classe} & \multicolumn{3}{c|}{2ª classe} & \multicolumn{3}{c|}{1ª classe} \\ \hline
\multicolumn{3}{|c|}{milhões} & \multicolumn{3}{c|}{milhares} & \multicolumn{3}{c|}{-} \\ \hline
\multicolumn{1}{|c|}{9ª O} & \multicolumn{1}{c|}{8ª O} & 7ª O & \multicolumn{1}{c|}{6ª O} & \multicolumn{1}{c|}{5ª O} & 4ª O & \multicolumn{1}{c|}{3ª O} & \multicolumn{1}{c|}{2ª O} & 1ª O \\ \hline
\multicolumn{1}{|c|}{C} & \multicolumn{1}{c|}{D} & U & \multicolumn{1}{c|}{C} & \multicolumn{1}{c|}{D} & U & \multicolumn{1}{c|}{C} & \multicolumn{1}{c|}{D} & U \\ \hline
\rowcolor[HTML]{FFFC9E} 
\multicolumn{1}{|c|}{\cellcolor[HTML]{FFFC9E}\textbf{}} & \multicolumn{1}{c|}{\cellcolor[HTML]{FFFC9E}\textbf{4}} & \textbf{5} & \multicolumn{1}{c|}{\cellcolor[HTML]{FFFC9E}\textbf{8}} & \multicolumn{1}{c|}{\cellcolor[HTML]{FFFC9E}\textbf{7}} & \textbf{4} & \multicolumn{1}{c|}{\cellcolor[HTML]{FFFC9E}\textbf{6}} & \multicolumn{1}{c|}{\cellcolor[HTML]{FFFC9E}\textbf{7}} & \textbf{4} \\ \hline
\end{tabular}
\end{table}

    \item 854.460.705 = 800.000.000 + 50.000.000 + 4.000.000 + 400.000 + 60.000 + 700 + 5

\begin{table}[H]
\begin{tabular}{|ccc|ccc|ccc|}
\hline
\multicolumn{3}{|c|}{3ª classe} & \multicolumn{3}{c|}{2ª classe} & \multicolumn{3}{c|}{1ª classe} \\ \hline
\multicolumn{3}{|c|}{milhões} & \multicolumn{3}{c|}{milhares} & \multicolumn{3}{c|}{-} \\ \hline
\multicolumn{1}{|c|}{9ª O} & \multicolumn{1}{c|}{8ª O} & 7ª O & \multicolumn{1}{c|}{6ª O} & \multicolumn{1}{c|}{5ª O} & 4ª O & \multicolumn{1}{c|}{3ª O} & \multicolumn{1}{c|}{2ª O} & 1ª O \\ \hline
\multicolumn{1}{|c|}{C} & \multicolumn{1}{c|}{D} & U & \multicolumn{1}{c|}{C} & \multicolumn{1}{c|}{D} & U & \multicolumn{1}{c|}{C} & \multicolumn{1}{c|}{D} & U \\ \hline
\rowcolor[HTML]{FFFC9E} 
\multicolumn{1}{|c|}{\cellcolor[HTML]{FFFC9E}\textbf{8}} & \multicolumn{1}{c|}{\cellcolor[HTML]{FFFC9E}\textbf{5}} & \textbf{4} & \multicolumn{1}{c|}{\cellcolor[HTML]{FFFC9E}\textbf{4}} & \multicolumn{1}{c|}{\cellcolor[HTML]{FFFC9E}\textbf{6}} & \textbf{0} & \multicolumn{1}{c|}{\cellcolor[HTML]{FFFC9E}\textbf{7}} & \multicolumn{1}{c|}{\cellcolor[HTML]{FFFC9E}\textbf{0}} & \textbf{5} \\ \hline
\end{tabular}
\end{table}

\end{itemize}

E dessa forma podemos ter um número infinito de exemplos, aumentando o número de classes sempre que for preciso (4ª bilhão, 5ª trilhão etc).

O objetivo deste capítulo não foi preparar você, aluno(a), para colocar números em tabelas separando por unidades, dezenas e centenas em suas respectivas classes, mas sim para apresentar o conceito de decomposição dos números em “\textbf{algarismos}” (os números que colocamos dentro de cada ordem) que será um conhecimento muito útil futuramente no estudo de notação científica e até na compreensão de conceitos mais simples como soma, subtração, multiplicação e divisão.